% =====================================================================
%  CARNET D'ENTRAINEMENT — FICHE 6 — THÈME 2
%  Niveau FACILE — Corrigé
% =====================================================================
\documentclass[11pt,a4paper]{article}
\usepackage[utf8]{inputenc}
\usepackage[T1]{fontenc}
\usepackage{lmodern}
\usepackage[french]{babel}
\usepackage{amsmath,amssymb,mathtools}
\usepackage{icomma}
\usepackage[version=4]{mhchem}
\usepackage{siunitx}
\sisetup{output-decimal-marker={,}}
\usepackage{xcolor}
\usepackage{colortbl}
\usepackage{tikz}
\usepackage[most]{tcolorbox}
\usepackage{enumitem}
\usepackage{array}
\usepackage{booktabs}
\usepackage{fancyhdr}
\usepackage[hidelinks]{hyperref}
\usetikzlibrary{arrows.meta,positioning,calc,shapes.geometric}
\usepackage[a4paper,margin=2.1cm,headheight=26pt,headsep=8pt]{geometry}

\definecolor{primary}{RGB}{146,95,160}
\definecolor{primarydark}{RGB}{92,52,104}
\definecolor{corcol}{RGB}{0,114,178}
\definecolor{astucecol}{RGB}{198,93,9}

\tcbset{boxbase/.style={enhanced,breakable,boxrule=0.9pt,arc=2.5pt,
   left=3mm,right=3mm,top=2.2mm,bottom=2.2mm,
   fonttitle=\bfseries,coltitle=white,toptitle=0.6mm,bottomtitle=0.6mm}}
\newtcolorbox[auto counter]{cor}[1][]{boxbase,colback=corcol!5,colframe=corcol,
   title={\textsc{Corrigé}~\thetcbcounter\ifx\relax#1\relax\relax\else\textmd{ —\itshape\ #1}\fi}}
\newtcolorbox{astucebox}[1][]{boxbase,colback=astucecol!8,colframe=astucecol,
   title={\textsc{Astuce}\ifx\relax#1\relax\relax\else\textmd{ : \itshape #1}\fi}}

% -------- macros des quatre grandeurs --------
\newcommand{\nmax}{n_{\max}(e^-)}
\newcommand{\nv}{n_{v}(e^-)}
\newcommand{\nl}{n_{\ell}(e^-)}
\newcommand{\nnl}{n_{n\ell}(e^-)}

\pagestyle{fancy}\fancyhf{}
\fancyhead[L]{\small\textcolor{primarydark}{\textbf{Fiche 6 · Thème 2} — Méthode des 5 étapes et structure de Lewis}}
\fancyhead[R]{\small\textcolor{primarydark}{\textsc{Facile · Corrigé}}}
\fancyfoot[C]{\small\thepage}
\renewcommand{\headrulewidth}{0.8pt}
\renewcommand{\headrule}{\hbox to\headwidth{\color{primary}\leaders\hrule height \headrulewidth\hfill}}
\setlength{\parindent}{0pt}\setlength{\parskip}{3pt}

\newcommand{\rep}{\textbf{\textcolor{corcol}{Réponse : }}}

\begin{document}
\begin{center}
{\Large\bfseries\color{primarydark}Thème 2 — Corrigé du niveau Facile}
\end{center}
\medskip

\begin{cor}[calculer $\nmax = 8a + 2b$]
On compte $a$ = atomes visant l'octet (tous sauf \ce{H}) et $b$ = nombre de \ce{H}.
\begin{enumerate}[label=\alph*),leftmargin=2em,itemsep=2pt]
  \item \ce{NH3} : $a=1$ (\ce{N}), $b=3$ (les \ce{H}) $\Rightarrow \nmax = 8\times 1 + 2\times 3 = \mathbf{14}$.
  \item \ce{H2O} : $a=1$ (\ce{O}), $b=2$ $\Rightarrow \nmax = 8 + 4 = \mathbf{12}$.
  \item \ce{CH4} : $a=1$ (\ce{C}), $b=4$ $\Rightarrow \nmax = 8 + 8 = \mathbf{16}$.
  \item \ce{CO2} : $a=3$ (\ce{C} et $2\,\ce{O}$), $b=0$ $\Rightarrow \nmax = 8\times 3 + 0 = \mathbf{24}$.
\end{enumerate}
\rep $14$ ; $12$ ; $16$ ; $24$.
\end{cor}

\begin{cor}[calculer $\nv$ pour des composés neutres]
Édifices neutres : $q=0$, donc $\nv$ est la simple somme des valences.
\begin{enumerate}[label=\alph*),leftmargin=2em,itemsep=2pt]
  \item \ce{NH3} : $\nv = 5 + 3\times 1 = \mathbf{8}$.
  \item \ce{H2O} : $\nv = 6 + 2\times 1 = \mathbf{8}$.
  \item \ce{CCl4} : $\nv = 4 + 4\times 7 = \mathbf{32}$.
  \item \ce{CO2} : $\nv = 4 + 2\times 6 = \mathbf{16}$.
\end{enumerate}
\rep $8$ ; $8$ ; $32$ ; $16$.
\end{cor}

\begin{cor}[du couple $(\nmax\,;\,\nv)$ au nombre de liaisons]
On applique $\nl = \nmax - \nv$, puis on divise par $2$.
\begin{enumerate}[label=\alph*),leftmargin=2em,itemsep=2pt]
  \item \ce{NH3} : $\nl = 14 - 8 = 6 \Rightarrow 6/2 = \mathbf{3}$ liaisons.
  \item \ce{CO2} : $\nl = 24 - 16 = 8 \Rightarrow 8/2 = \mathbf{4}$ liaisons (les deux doubles liaisons \ce{C=O}).
  \item \ce{H2O} : $\nl = 12 - 8 = 4 \Rightarrow 4/2 = \mathbf{2}$ liaisons.
\end{enumerate}
\rep $3$ liaisons ; $4$ liaisons ; $2$ liaisons.
\end{cor}

\begin{cor}[du couple $(\nv\,;\,\nl)$ au nombre de doublets libres]
On applique $\nnl = \nv - \nl$, puis on divise par $2$.
\begin{enumerate}[label=\alph*),leftmargin=2em,itemsep=2pt]
  \item \ce{NH3} : $\nnl = 8 - 6 = 2 \Rightarrow 2/2 = \mathbf{1}$ doublet libre.
  \item \ce{CO2} : $\nnl = 16 - 8 = 8 \Rightarrow 8/2 = \mathbf{4}$ doublets libres.
  \item \ce{H2O} : $\nnl = 8 - 4 = 4 \Rightarrow 4/2 = \mathbf{2}$ doublets libres.
\end{enumerate}
\rep $1$ doublet ; $4$ doublets ; $2$ doublets.
\end{cor}

\begin{cor}[calculer $\nv$ pour des ions : gare au $-q$]
On soustrait la charge : $-q$ \emph{ajoute} des électrons à un anion, en \emph{retranche} à un
cation.
\begin{enumerate}[label=\alph*),leftmargin=2em,itemsep=2pt]
  \item \ce{OH^-} ($q=-1$) : $\nv = 6 + 1 - (-1) = 6 + 1 + 1 = \mathbf{8}$.
  \item \ce{NH4+} ($q=+1$) : $\nv = 5 + 4\times 1 - (+1) = 5 + 4 - 1 = \mathbf{8}$.
  \item \ce{CO3^2-} ($q=-2$) : $\nv = 4 + 3\times 6 - (-2) = 4 + 18 + 2 = \mathbf{24}$.
  \item \ce{NO3^-} ($q=-1$) : $\nv = 5 + 3\times 6 - (-1) = 5 + 18 + 1 = \mathbf{24}$.
\end{enumerate}
\rep $8$ ; $8$ ; $24$ ; $24$.
\end{cor}

\begin{astucebox}[le bon signe du premier coup]
Pour un \textbf{anion} de charge $-n$, on \emph{ajoute} $n$ (l'ion a capté $n$ électrons
supplémentaires) ; pour un \textbf{cation} $+n$, on \emph{retranche} $n$. Puisque
$\nv = \sum N_v - q$, il suffit de reporter $q$ avec son signe : $-(-2) = +2$ pour \ce{CO3^2-},
$-(+1) = -1$ pour \ce{NH4+}.
\end{astucebox}

\end{document}
